\section{A worked example}\label{sec:worked}

Four data points are enough to show the whole construction, and to exhibit in arithmetic
the main properties Section~\ref{sec:method} proves in general. Take one cell containing
four test cases and two possible labels. Three of the cases carry label $1$ and one
carries label $2$, so the cell's label frequencies are $(3/4,1/4)$. The old model $q_0$
is uninformative within the cell: it predicts $(0.5,0.5)$ for all four cases. We compare
two candidate replacements, shown in Table~\ref{tab:worked}.

\begin{table}[t]
\centering
\small
\caption{A four-point example in a single cell. Model A moves every case to the cell's
label frequencies; model B moves each case towards its own correct label by the same
amount. The two models are equally far from $q_0$, and only B uses information about the
individual case.}
\label{tab:worked}
\begin{tabular}{@{}ccccc@{}}
\toprule
case $i$ & label $y_i$ & $q_0(\cdot\mid x_i)$ & model A: $q_1(\cdot\mid x_i)$ &
model B: $q_1(\cdot\mid x_i)$ \\
\midrule
1 & 1 & $(0.5,0.5)$ & $(0.75,0.25)$ & $(0.75,0.25)$ \\
2 & 1 & $(0.5,0.5)$ & $(0.75,0.25)$ & $(0.75,0.25)$ \\
3 & 1 & $(0.5,0.5)$ & $(0.75,0.25)$ & $(0.75,0.25)$ \\
4 & 2 & $(0.5,0.5)$ & $(0.75,0.25)$ & $(0.25,0.75)$ \\
\bottomrule
\end{tabular}
\end{table}

Score both models with the quadratic score $S(q,y)=2q_y-\lVert q\rVert_2^2$, and for each
case tabulate the score \emph{difference} at each of the two possible labels, not only at
the observed one. Writing $h_i(y)=S(q_1(\cdot\mid x_i),y)-S(q_0(\cdot\mid x_i),y)$, all
the arithmetic is in eighths:
\begin{equation}
h_i(y)=
\begin{cases}
+3/8 & \text{if $y$ is the label the new model moved towards},\\
-5/8 & \text{otherwise}.
\end{cases}
\label{eq:worked-h}
\end{equation}
This vector, evaluable at labels that were never observed for case $i$, is the only object
the method needs beyond the labels themselves.

\paragraph{Model A: a pure prior refit.}
Every case gets the same prediction, so $h_i(1)=+3/8$ and $h_i(2)=-5/8$ for all four
cases. The \emph{observed} mean gain uses each case's own label,
\begin{equation}
\Delta T=\tfrac14\left[3\cdot\tfrac38+1\cdot\left(-\tfrac58\right)\right]=\tfrac18 .
\end{equation}
Now transport labels: replace each case's label, in turn, by each of the other three cases'
labels and average. Case~1 is scored against labels $\{1,1,2\}$, giving
$\frac13(\frac38+\frac38-\frac58)=\frac{1}{24}$, and likewise for cases 2 and 3; case~4 is
scored against $\{1,1,1\}$, giving $+\frac38$. The mean transported gain is
\begin{equation}
\Delta P=\tfrac14\left[3\cdot\tfrac{1}{24}+\tfrac38\right]=\tfrac18 ,
\qquad\text{so}\qquad
\Delta R=\Delta T-\Delta P=0 .
\end{equation}
Model A's entire advantage survives the relabelling, and its covariance gain is exactly
zero --- not approximately, and not zero on average over hypothetical resamples, but zero
in this sample. That is the general fact: a prediction change that is constant within a
cell contributes nothing to $\Delta R$ (Theorem~\ref{thm:population}). Note also that
$\Delta P=1/8$ is the largest transported gain any within-cell-constant prediction can earn
here, because a proper score is maximized in expectation at the true frequencies
$(3/4,1/4)$ --- which is exactly what model A predicts.

\paragraph{Model B: a case-specific improvement.}
Cases 1--3 are unchanged, but case~4 now moves towards label $2$, so
$h_4(1)=-5/8$ and $h_4(2)=+3/8$. Every case is now scored at $+3/8$ on its own label, so
$\Delta T=3/8$. Under transport, cases 1--3 are unaffected ($\frac{1}{24}$ each), but
case~4 is scored against three label-$1$ cases at $-5/8$ apiece. The mean transported gain
turns negative,
\begin{equation}
\Delta P=\tfrac14\left[3\cdot\tfrac{1}{24}-\tfrac58\right]=-\tfrac18 ,
\qquad
\Delta R=\tfrac38-\left(-\tfrac18\right)=\tfrac12 .
\end{equation}
Two things are worth reading off this. First, model B's aggregate gain of $3/8$
\emph{understates} its case-level improvement of $1/2$, and the reason is \emph{not} that
B fits the cell's label frequencies worse. It fits them better: averaged over the cell it
predicts $0.625$ for label 1 against the old model's $0.5$, and the truth is $0.75$, so
its mean-forecast fit improves by exactly $0.09375$. What offsets that is dispersion.
B's forecasts vary inside the cell where the old model's did not, and by
Proposition~\ref{prop:transported} the transported channel charges that variation at
exactly the same $0.09375$; the two brackets cancel, so the transported gain computed with
the cell's own frequencies is exactly zero. The $-1/8$ the estimator reports is what
remains after each case is scored against the \emph{other} cases' labels only, which
removes the $\tfrac1{n_c}\widehat{\Delta R}=1/8$ that the in-sample plug-in had credited
to the transported channel. Both numbers are right; they answer slightly different
questions, and Proposition~\ref{prop:attenuate} relates them exactly.

Second, an evaluator reading only the total would credit model B with three times model
A's improvement, whereas the decomposition shows the two improving on disjoint grounds: A
earns $1/8$ purely by moving to the cell's frequencies and gains no covariance at all,
while B's $1/2$ of covariance gain is partly concealed by a transported term that turns
negative. This is the failure mode that motivates the paper, in miniature: aggregate
scores are not merely imprecise about composition, they can order two models by the wrong
criterion.

\paragraph{The covariance form, and why leaving one out matters.}
Theorem~\ref{thm:population} says $\Delta R$ is a within-cell covariance between the
probability change and the label indicator. Computing that covariance directly for model
B, over the four cases, gives
$2\sum_y\operatorname{Cov}(\Delta q_y,\ind\{Y=y\})=3/8$ --- not the $1/2$ the transport
calculation returned. The discrepancy is not an error in either: the covariance uses the
cell's \emph{own} four labels to form the reference frequencies, so each case contributes
to the reference it is compared against. That in-sample plug-in is attenuated by exactly
$(n_c-1)/n_c$, here $3/4$, and indeed $\frac34\cdot\frac12=\frac38$. Transport, which
scores each case only against the \emph{other} cases' labels, removes the attenuation
without any held-out data. Proposition~\ref{prop:attenuate} shows this holds exactly, for
every cell with $n_c\ge2$ and every grouping variable, which is why the method needs no
sample splitting. The example is reproduced as a unit test in the accompanying software.

One property the example cannot show, because both of its candidate models happen to be
as confident as each other, is the estimator's sensitivity to confidence itself. Multiply
every prediction's odds by a constant and $\Delta R$ moves, even though no case has been
reordered and no information added. That is a real limitation rather than a technicality
--- Section~\ref{sec:simulation} measures it at $-0.488$ in a design where the two models
carry identical case-level information --- and it is why every result in this paper is
reported for confidence-matched pairs.

\subsection{A decision the split changes}\label{sec:decision}

A decomposition that no decision depends on would be an ornament. The four points above
are enough to show one, and to show that the split predicts it where the aggregate score
cannot.

Read at the benchmark's own label frequencies, both candidate models improve on $q_0$ and
model B improves by three times as much: $\Delta T=1/8$ against $3/8$. An evaluator
choosing between them on the reported score picks B, and is right to. Now suppose the
system will be deployed where the label frequencies in this cell are
$(\tfrac14,\tfrac34)$ rather than the $(\tfrac34,\tfrac14)$ the test set happens to
exhibit --- a label shift, the situation post-hoc prior correction exists to address
\citep{saerens2002adjusting,lipton2018detecting}. Reweighting each case by
$p'(y_i)/p(y_i)$ and rescoring gives
\begin{equation}
\text{model A}:\ +\tfrac18\ \longrightarrow\ -\tfrac38,
\qquad
\text{model B}:\ +\tfrac38\ \longrightarrow\ +\tfrac38 .
\end{equation}
Model A does not merely lose its margin; it becomes \emph{worse} than the model it was
meant to replace, while model B is untouched. The ranking of A against $q_0$ reverses on a
change in the deployment distribution that leaves every model and every prediction
exactly as it was.

The decomposition says in advance which of the two is exposed, and the aggregate score
does not. A's entire gain sits in the transported channel, which is the channel a change
in label frequencies acts on; B's sits in the covariance channel, which such a change
leaves alone, because scoring each case against its own label is unaffected by how often
that label occurs once the case is reweighted. This is the operational content of the
split: $\widehat{\Delta P}$ is the part of a reported improvement that is contingent on
the evaluation set's label frequencies matching the deployment's, and
$\widehat{\Delta R}$ is the part that is not. Two systems with the same headline gain and
different compositions are different bets, and only the decomposition prices them
differently. Section~\ref{sec:discussion} returns to what this does and does not license,
and Appendix~\ref{app:vgvisual} exhibits the same structure on real models, where
correcting a model's within-cell prior post hoc registers as almost pure transported gain
($+0.02401$ against $+0.00111$) exactly as a correction carrying no case-level
information must.
